\documentclass[a4paper,ngerman,14pt,]{scrlttr2}

\usepackage[margin=20mm,footskip=10mm]{geometry}
\usepackage[utf8]{inputenc}

\usepackage[nolinks]{qrcode}
\usepackage{graphicx}
\usepackage{tikz}

\usepackage[light,semibold]{noto}
\usepackage[T1]{fontenc}

\begin{document}


{% for entry in data.entries %}

\thispagestyle{empty}
\begin{tikzpicture}[remember picture,overlay]
   \node[anchor=north east,inner sep=5pt] at (current page.north east)
   {\includegraphics[scale=0.6]{rwth_asta_cmyk.eps} };
\end{tikzpicture}
\begin{center}
{\huge{\textbf{Ticket Check-In ({{ data.event }}) } } }\\
\vspace{3mm}
1. Du benötigst dein (geladenes) Smartphone und die App \textit{PretixSCAN}.\\
\vspace{3mm}
Lade die App für Android oder iOS herunter:\\
\vspace{3mm}
\begin{tabular}{lr}
\qrcode[height=0.4\linewidth]{https://play.google.com/store/apps/details?id=eu.pretix.pretixscan.droid}
&
\hspace{0.1\linewidth}
\qrcode[height=0.4\linewidth]{https://apps.apple.com/de/app/pretixscan/id1465873864}
\\
\tiny{https://play.google.com/store/apps/details?id=eu.pretix.pretixscan.droid}
&
\tiny{https://apps.apple.com/de/app/pretixscan/id1465873864}
\end{tabular}\\
\vspace{5mm}
2. Öffne die App und setze sie ggf. auf den Werkszustand zurück.\\
\vspace{3mm}
3. Scanne den folgenden Konfigurations-Code:\\
\vspace{8mm}
{{ entry.qrdata }}\\
\vspace{1mm}
%{\footnotesize
{\tiny{
Wenn die App oder das Gerät das Scannen von QR-Codes nicht unterstützt, können Sie auch die folgenden Daten manuell eingeben:\\
System-Adresse: \texttt{ {{ data.url }} }\\
Token: \texttt{ {{ entry.token }} }\\
}}
\vspace{5mm}
ID: {{ entry.name }}
\end{center}
\newpage

{% endfor %}


\end{document}
